\documentclass{BeamerTemplate}

\title{Module 8 - Analyse de la variance (ANOVA)}

\subtitle{STT-1920 Méthodes statistiques}
\date{\today}

\begin{document}
\begin{frame}[t]
  \titlepage
\end{frame}

%%%%%%%%%%%%%%%%%%%%%%%%%%%%%%%%%%%
\section{ANOVA à un facteur}
%%%%%%%%%%%%%%%%%%%%%%%%%%%%%%%%%%%

\begin{frame}[t]{Motivation}

On veut comparer les moyennes de \alert{$I \ge 3$ groupes} (traitements, niveaux d'un facteur).

\bigskip
\underline{Exemple :} Comparer la distance de freinage de 3 marques de freins à disque ($I=3$), avec $k$ mesures par marque.

\bigskip
\begin{itemize}\itemsep4mm
  \item Effectuer plusieurs tests $t$ deux à deux multiplie les erreurs de type I.
  \item L'ANOVA teste l'égalité de toutes les moyennes en une seule procédure.
\end{itemize}

\bigskip
$H_0 : \mu_1 = \mu_2 = \cdots = \mu_I \quad$ contre $\quad H_1 :$ au moins une $\mu_i$ diffère.

\end{frame}

\begin{frame}[t]{Modèle d'ANOVA à un facteur}

\begin{Boite}{Modèle}
$$Y_{ij} = \mu + \tau_i + \varepsilon_{ij}, \quad i=1,\ldots,I,\; j=1,\ldots,n_i$$
\end{Boite}

\begin{itemize}\itemsep4mm
  \item $\mu$ : moyenne générale.
  \item $\tau_i$ : effet du $i$\textsuperscript{e} traitement, avec $\sum_i \tau_i = 0$.
  \item $\varepsilon_{ij} \sim \mathcal{N}(0, \sigma^2)$ i.i.d. : erreur aléatoire.
\end{itemize}

\bigskip
\begin{Boite}{Postulats du modèle}
\begin{itemize}
  \item Normalité des erreurs.
  \item Homogénéité des variances ($\sigma^2$ identique dans tous les groupes).
  \item Indépendance des observations.
\end{itemize}
\end{Boite}

\end{frame}

\begin{frame}[t]{Décomposition de la variance}

L'idée clé : décomposer la variabilité totale en variabilité \alert{entre groupes} et variabilité \alert{au sein des groupes}.

\bigskip
$$SS_T = SS_{\text{Trait}} + SS_E$$

\medskip
\begin{itemize}\itemsep4mm
  \item $SS_T = \sum_{i,j}(Y_{ij} - \overline{Y})^2$ : variabilité totale.
  \item $SS_{\text{Trait}} = \sum_i n_i(\overline{Y}_{i\cdot} - \overline{Y})^2$ : variabilité entre groupes.
  \item $SS_E = \sum_{i,j}(Y_{ij} - \overline{Y}_{i\cdot})^2$ : variabilité au sein des groupes (erreur).
\end{itemize}

\end{frame}

\begin{frame}[t]{Table d'ANOVA}

\begin{center}
\begin{tabular}{lccccc}\hline
\textbf{Source} & \textbf{SS} & \textbf{ddl} & \textbf{MS} & \textbf{$F_0$} & \textbf{Valeur-$p$} \\ \hline
Traitement & $SS_{\text{Trait}}$ & $I-1$ & $MS_{\text{Trait}}$ & $F_0 = MS_{\text{Trait}}/MS_E$ & \\[3pt]
Erreur & $SS_E$ & $n-I$ & $MS_E$ & & \\[3pt]
Total & $SS_T$ & $n-1$ & & & \\ \hline
\end{tabular}
\end{center}

\medskip
\begin{itemize}\itemsep3mm
  \item $MS = SS / \text{ddl}$ : carré moyen.
  \item $MS_E = \hat{\sigma}^2$ : estimateur de la variance.
  \item Sous $H_0$ : $F_0 \sim F_{I-1,\,n-I}$.
  \item Rejeter $H_0$ si $F_0 > F_{\alpha;\,I-1,\,n-I}$ ou si valeur-$p < \alpha$.
\end{itemize}

\end{frame}

\begin{frame}[t]{Exemple — freins à disque}

\begin{verbatim}
                   SS   DDL       MS       F0    valeur-p
Traitement      45.77     2   22.885   14.329    0.000001
Erreur          91.03    57    1.597
\end{verbatim}

\begin{itemize}\itemsep4mm
  \item $I-1=2 \Rightarrow I=3$ marques de freins.
  \item $n-I=57 \Rightarrow n=60$ observations. Donc $k=20$ répétitions par marque.
  \item $\hat{\sigma}^2 = MS_E = 1{,}597$.
  \item $F_0 = 14{,}329$, valeur-$p = 0{,}000001 < 0{,}01 = \alpha$.
  \item \alert{On rejette $H_0$} : au moins une moyenne de distance de freinage diffère.
\end{itemize}

\end{frame}

%%%%%%%%%%%%%%%%%%%%%%%%%%%%%%%%%%%
\section{ANOVA à deux facteurs}
%%%%%%%%%%%%%%%%%%%%%%%%%%%%%%%%%%%

\begin{frame}[t]{ANOVA à deux facteurs}

On s'intéresse à l'effet de \alert{deux facteurs} A et B sur la variable réponse.

\bigskip
\begin{Boite}{Modèle avec interaction}
$$Y_{ijk} = \mu + \tau_i + \beta_j + (\tau\beta)_{ij} + \varepsilon_{ijk}$$
\end{Boite}

\begin{itemize}\itemsep4mm
  \item $\tau_i$ : effet principal du facteur A.
  \item $\beta_j$ : effet principal du facteur B.
  \item $(\tau\beta)_{ij}$ : \alert{terme d'interaction} — l'effet de A dépend du niveau de B.
\end{itemize}

\end{frame}

\begin{frame}[t]{Interaction}

\begin{Boite}{Interaction}
L'effet d'un facteur dépend du niveau de l'autre facteur.
\end{Boite}

\medskip
\begin{itemize}\itemsep5mm
  \item \textbf{Pas d'interaction} : les courbes des profils moyens sont parallèles.
  \item \textbf{Interaction présente} : les courbes se croisent.
\end{itemize}

\bigskip
Si l'interaction est significative, on ne peut pas interpréter les effets principaux indépendamment.

\end{frame}

\begin{frame}[t]{Table d'ANOVA à deux facteurs}

\begin{center}
\small
\begin{tabular}{lcccc}\hline
\textbf{Source} & \textbf{ddl} & \textbf{MS} & \textbf{$F_0$} \\ \hline
Facteur A & $I-1$ & $MS_A$ & $MS_A/MS_E$ \\[3pt]
Facteur B & $J-1$ & $MS_B$ & $MS_B/MS_E$ \\[3pt]
Interaction $A \times B$ & $(I-1)(J-1)$ & $MS_{AB}$ & $MS_{AB}/MS_E$ \\[3pt]
Erreur & $IJ(n-1)$ & $MS_E$ & \\[3pt]
Total & $IJn-1$ & & \\ \hline
\end{tabular}
\end{center}

\medskip
On teste d'abord l'interaction, puis les effets principaux si pertinent.

\end{frame}

\begin{frame}[t]{Comparaisons multiples (post-hoc)}

Après un test $F$ significatif, on veut savoir \alert{quels groupes diffèrent}.

\bigskip
\begin{Boite}{Méthode de Tukey (HSD)}
Contrôle le taux d'erreur familial à $\alpha$ pour l'ensemble des comparaisons deux à deux.
$$|\overline{y}_{i\cdot} - \overline{y}_{j\cdot}| > q_{\alpha;\,I,\,n-I}\,\frac{MS_E}{\sqrt{k}}$$
\end{Boite}

\medskip
D'autres méthodes existent : Bonferroni, Dunnett (comparaison à un contrôle), etc.

\end{frame}

\end{document}
